\documentclass[12pt,a4paper]{article}

% =========================================================
% Packages
% =========================================================
\usepackage[T1]{fontenc}
\usepackage[utf8]{inputenc}
\usepackage[top=2.5cm, bottom=2.5cm, left=3cm, right=2.5cm]{geometry}

\usepackage{graphicx}
\usepackage{amsmath}
\usepackage{booktabs}
\usepackage{array}
\usepackage{xcolor}
\usepackage{hyperref}
\usepackage{caption}
\usepackage{subcaption}
\usepackage{listings}
\usepackage{enumitem}
\usepackage{titlesec}
\usepackage{fancyhdr}
\usepackage{float}
\usepackage{parskip}
\usepackage{longtable}
\usepackage{multirow}

% =========================================================
% Colors
% =========================================================
\definecolor{primaryblue}{RGB}{31,78,121}
\definecolor{codeblue}{RGB}{0,114,178}
\definecolor{codegray}{RGB}{245,245,245}
\definecolor{commentgreen}{RGB}{40,160,40}
\definecolor{stringred}{RGB}{180,0,0}

% =========================================================
% Hyperref setup
% =========================================================
\hypersetup{
    colorlinks=true,
    linkcolor=primaryblue,
    citecolor=primaryblue,
    urlcolor=codeblue,
    pdftitle={CityPersons SD3.5 Pedestrian Augmentation: Technical Report},
    pdfauthor={Bui Duc Thang}
}

% =========================================================
% Code listing style
% =========================================================
\lstset{
    backgroundcolor=\color{codegray},
    basicstyle=\ttfamily\small,
    breaklines=true,
    commentstyle=\color{commentgreen},
    keywordstyle=\color{codeblue}\bfseries,
    stringstyle=\color{stringred},
    numberstyle=\tiny\color{gray},
    frame=single,
    framesep=4pt,
    rulecolor=\color{gray!40},
    showstringspaces=false,
    tabsize=4,
    captionpos=b,
    language=Python
}

% =========================================================
% Header / Footer
% =========================================================
\pagestyle{fancy}
\fancyhf{}
\fancyhead[L]{\small\textcolor{gray}{CityPersons SD3.5 Augmentation}}
\fancyhead[R]{\small\textcolor{gray}{Technical Report}}
\fancyfoot[C]{\thepage}
\renewcommand{\headrulewidth}{0.4pt}

% =========================================================
% Section formatting
% =========================================================
\titleformat{\section}
{\large\bfseries\color{primaryblue}}
{\thesection}{1em}{}
[\vspace{-4pt}\rule{\linewidth}{0.6pt}]

\titleformat{\subsection}
{\normalsize\bfseries\color{primaryblue!80!black}}
{\thesubsection}{1em}{}

\titleformat{\subsubsection}
{\normalsize\itshape\color{primaryblue!60!black}}
{\thesubsubsection}{1em}{}

% =========================================================
% Document
% =========================================================
\begin{document}

% =========================================================
% Title Page
% =========================================================
\begin{titlepage}
    \centering
    \vspace*{2cm}

    {\color{primaryblue}\rule{\linewidth}{2pt}}
    \vspace{0.5cm}

    {\Huge\bfseries\color{primaryblue}
        CityPersons Pedestrian Augmentation\\[6pt]
        using Stable Diffusion 3.5 Medium
    }

    \vspace{0.5cm}
    {\color{primaryblue}\rule{\linewidth}{2pt}}

    \vspace{1.5cm}

    {\Large\bfseries Technical Report}

    \vspace{1cm}

    \begin{tabular}{ll}
        \textbf{Project}    & CityPersons SD3.5 Pedestrian Augmentation \\[4pt]
        \textbf{Task}       & Synthetic pedestrian insertion for urban Camera AI \\[4pt]
        \textbf{Notebook}   & \texttt{v25 i think.ipynb} \\[4pt]
        \textbf{Main Model} & Stable Diffusion 3.5 Medium \\[4pt]
        \textbf{Segmentation Model} & YOLOv8m-seg \\[4pt]
        \textbf{Dataset}    & CityPersons \\[4pt]
        \textbf{Mode}       & Context person composite \\[4pt]
        \textbf{Hardware}   & Kaggle GPU environment \\[4pt]
        \textbf{Author}     & Bui Duc Thang \\
    \end{tabular}

    \vfill

    {\small\color{gray}
    VinSmart Future -- Camera AI / Computer Vision R\&D \quad|\quad June 2026
    }

\end{titlepage}

% =========================================================
% Table of Contents
% =========================================================
\tableofcontents
\clearpage

% =========================================================
% Abstract
% =========================================================
\begin{abstract}
\noindent
This report documents the current implementation of a synthetic pedestrian
augmentation pipeline for the CityPersons dataset using Stable Diffusion 3.5
Medium. The goal of the project is to insert realistic pedestrians into existing
urban street images while preserving the original background geometry, including
road layout, vehicles, buildings, camera viewpoint, and perspective structure.

The current notebook no longer follows a LoRA-training-first workflow. Instead,
it implements a practical generation pipeline based on Smart Placement,
context-aware SD3.5 image-to-image generation, YOLOv8m-seg person segmentation,
person-only compositing, color matching, retry logic, debug export, and manifest
logging. The active background preservation mode is
\texttt{context\_person\_composite}. In this mode, Stable Diffusion generates a
candidate pedestrian inside a local context crop, YOLOv8m-seg extracts the
generated person mask, and only the pedestrian pixels are pasted back onto the
original image.

The main technical insight is that background preservation should not be solved
only through prompt engineering. Full-image diffusion generation gives the model
too much freedom and can rewrite the scene. Therefore, the notebook separates
object generation from scene preservation: Stable Diffusion is used as a
candidate pedestrian generator, while the original image remains the source of
truth for the background.



\end{abstract}

\clearpage

% =========================================================
% 1. Introduction
% =========================================================
\section{Introduction}

Pedestrian perception is a core component of Camera AI systems. In urban
driving environments, systems for traffic monitoring, autonomous perception,
smart surveillance, and pedestrian safety rely on robust detection and tracking
of people at different distances, scales, occlusion levels, and lighting
conditions. Although modern object detectors have improved significantly, their
performance still depends heavily on the diversity and quality of annotated
training data.

CityPersons provides realistic urban pedestrian scenes, but the distribution of
pedestrian appearances and placements is still limited by the original collected
data. Synthetic augmentation can increase the diversity of the dataset by adding
new pedestrians into existing scenes. However, this task is difficult because
the inserted pedestrian must be realistic and must not corrupt the original
scene.

This project studies how Stable Diffusion 3.5 Medium can be used to generate
synthetic pedestrians for CityPersons. The main challenge is not simply to
generate a person. The generated person must be placed at a plausible position,
match perspective, have a realistic scale, touch the ground, and integrate with
the scene. At the same time, the background must remain stable so that road
geometry, vehicles, buildings, and camera viewpoint are preserved.

\subsection{Project Objective}

The objective of the notebook is:

\begin{quote}
To generate realistic synthetic pedestrians inside CityPersons urban street
images using Stable Diffusion 3.5 Medium, while preserving original background
geometry through YOLOv8m-seg based person-only compositing.
\end{quote}

The output images are intended to support downstream pedestrian detection and
tracking experiments. Therefore, visual realism alone is not enough. The
generated pedestrian should also be detectable by a person detector and suitable
for use as training data.

\subsection{Current Notebook Direction}

The current notebook is not mainly a LoRA training notebook. Earlier project
iterations considered LoRA fine-tuning, but the current implementation focuses
on a no-LoRA baseline using:

\begin{itemize}
    \item Smart Placement for insertion-region selection;
    \item Stable Diffusion 3.5 Medium for context-aware image-to-image generation;
    \item YOLOv8m-seg for generated-person segmentation;
    \item person-only compositing for background preservation;
    \item color matching for scene consistency;
    \item retry logic for failed generations;
    \item debug image export and manifest logging.
\end{itemize}

LoRA remains a possible future direction, but only if the current baseline fails
systematically after parameter tuning, prompt engineering, placement adjustment,
and compositing refinement.

% =========================================================
% 2. Problem Definition
% =========================================================
\section{Problem Definition}

The task is synthetic pedestrian insertion. Given an original CityPersons image
$I$, the notebook generates an augmented image $I'$ such that:

\begin{enumerate}
    \item $I'$ contains at least one additional pedestrian.
    \item The new pedestrian is realistic and visible.
    \item The new pedestrian is placed at a plausible ground location.
    \item The pedestrian scale follows scene perspective.
    \item The original background is preserved outside the inserted person.
    \item The output image remains useful for detector or tracker training.
\end{enumerate}

This can be written as:

\begin{equation}
I' = C(I, G(I_{crop}, M, p), S)
\end{equation}

where:

\begin{itemize}
    \item $I$ is the original image;
    \item $I_{crop}$ is a local context crop around the insertion region;
    \item $M$ is the inpainting or generation mask;
    \item $p$ is the text prompt;
    \item $G$ is Stable Diffusion 3.5 generation;
    \item $S$ is the YOLOv8m-seg person mask;
    \item $C$ is the compositing function.
\end{itemize}

The key point is that the final background is primarily taken from $I$, not from
the generated crop. The generated crop is used to obtain the new pedestrian.

% =========================================================
% 3. Dataset
% =========================================================
\section{Dataset}

\subsection{CityPersons}

The primary dataset is CityPersons, an urban pedestrian detection dataset
derived from Cityscapes. It provides street-scene images with pedestrian
annotations. The dataset is suitable for this project because it contains urban
roads, sidewalks, vehicles, buildings, and pedestrians with realistic
perspective.

\begin{table}[H]
\centering
\caption{CityPersons dataset summary.}
\label{tab:citypersons_summary}
\begin{tabular}{lll}
\toprule
\textbf{Property} & \textbf{Value} & \textbf{Relevance} \\
\midrule
Training images & 2975 & Main source images for augmentation \\
Validation images & 500 & Validation / review split \\
Test images & 1575 & Reserved evaluation split \\
Average pedestrians per image & Around 7 & Useful pedestrian context \\
Annotations & Visible and full-body boxes & Support placement and scale reasoning \\
Scene type & Urban street scenes & Relevant to Camera AI \\
\bottomrule
\end{tabular}
\end{table}

\subsection{Why CityPersons Fits the Task}

CityPersons is suitable because the images contain stable urban structure and
realistic pedestrian scale. This is important for synthetic insertion because
the model needs contextual cues such as road surface, sidewalk, vehicles, and
camera perspective.

Compared with highly crowded datasets, CityPersons is easier to control for
insertion. Compared with tracking datasets such as MOT17, CityPersons provides
more static urban-scene diversity for augmentation. MOT17 can still be used
later for downstream tracking evaluation, but CityPersons is the main generation
dataset in the current notebook.

% =========================================================
% 4. Model Choice
% =========================================================
\section{Model Choice}

\subsection{Stable Diffusion 3.5 Medium}

The notebook uses Stable Diffusion 3.5 Medium as the main generative backbone.
The model is selected because it provides a practical balance between generation
quality and feasibility on Kaggle GPU environments.

The notebook uses:

\begin{lstlisting}[language=Python,caption={Current SD3.5 configuration}]
SD35_MODEL_ID = "stabilityai/stable-diffusion-3.5-medium"
RESOLUTION = 448
USE_T5 = False
\end{lstlisting}

The choice \texttt{USE\_T5=False} is important for Kaggle T4 safety. It reduces
memory requirements and makes the pipeline more practical in limited-VRAM
notebook environments.

\subsection{YOLOv8m-seg}

The current notebook uses YOLOv8m-seg for person segmentation:

\begin{lstlisting}[language=Python,caption={Current person segmentation model}]
CONTEXT_PERSON_SEGMENTATION_MODEL = "yolov8m-seg.pt"
\end{lstlisting}

This is an important detail. The notebook is not using \texttt{yolov8n-seg.pt}
in the current version. YOLOv8m-seg is heavier than YOLOv8n-seg but provides
stronger segmentation quality, which is useful because the final compositing
quality depends heavily on the person mask.

\subsection{Why Segmentation Is Needed}

Stable Diffusion generates a full crop, not only the person. If the generated
crop is directly pasted back into the original image, background pixels from the
generated crop may replace original pixels. This can create road changes,
sidewalk artifacts, building distortions, or inconsistent lighting.

YOLO segmentation solves this by extracting only the generated pedestrian. The
original image remains the background source.

% =========================================================
% 5. Current Notebook Configuration
% =========================================================
\section{Current Notebook Configuration}

The active configuration in \texttt{v25 i think.ipynb} is summarized below.

\begin{lstlisting}[language=Python,caption={Current global generation configuration}]
BACKGROUND_PRESERVATION_MODE = "context_person_composite"

RESOLUTION = 448
USE_T5 = False

AUGMENTATION_STRENGTH = 0.72
GUIDANCE_SCALE = 7.2
NUM_INFERENCE_STEPS = 42

CONTEXT_CROP_EXPAND = 3.4
CONTEXT_INPAINT_MASK_PADDING = 46
CONTEXT_GENERATION_RETRIES = 4

CONTEXT_PERSON_SEGMENTATION_MODEL = "yolov8m-seg.pt"
CONTEXT_PERSON_MIN_CONFIDENCE = 0.12
CONTEXT_PERSON_MASK_THRESHOLD = 0.32

COLOR_MATCH_PERSON_TO_SCENE = True

SAVE_PATCH_DEBUG = True
PATCH_DEBUG_MAX_ITEMS = 24

SCALE_ENVELOPE_HEIGHT_MULT = 1.16
SCALE_ENVELOPE_WIDTH_MULT = 1.30

PERSPECTIVE_SCALE_NEAR = 1.08
PERSPECTIVE_SCALE_FAR = 0.46
\end{lstlisting}

\subsection{Variant-Specific Parameters}

The notebook also defines variant-specific strength, guidance, and step values.

\begin{lstlisting}[language=Python,caption={Variant-specific strength values}]
VARIANT_STRENGTHS = {
    "add_single_pedestrian": 0.72,
    "add_two_pedestrians": 0.74,
    "add_small_group": 0.76,
    "add_occluded_pedestrian": 0.74,
    "add_distant_pedestrian": 0.68,
    "add_near_pedestrian": 0.76,
}
\end{lstlisting}

\begin{lstlisting}[language=Python,caption={Variant-specific guidance values}]
VARIANT_GUIDANCE_SCALES = {
    "add_single_pedestrian": 6.8,
    "add_two_pedestrians": 6.9,
    "add_small_group": 7.0,
    "add_occluded_pedestrian": 6.8,
    "add_distant_pedestrian": 6.6,
    "add_near_pedestrian": 7.3,
}
\end{lstlisting}

\begin{lstlisting}[language=Python,caption={Variant-specific inference steps}]
VARIANT_NUM_INFERENCE_STEPS = {
    "add_single_pedestrian": 38,
    "add_two_pedestrians": 38,
    "add_small_group": 40,
    "add_occluded_pedestrian": 38,
    "add_distant_pedestrian": 34,
    "add_near_pedestrian": 40,
}
\end{lstlisting}

\subsection{Current Smoke-Test Cell}

The current run cell generates 10 images using two variants:

\begin{lstlisting}[language=Python,caption={Current smoke-test cell}]
generated_paths = augment_dataset_limited_total(
    records,
    total_images=10,
    variants=["add_single_pedestrian", "add_near_pedestrian"],
    target_splits=["train"]
)

generated_paths[:5]
\end{lstlisting}

This means the current notebook is still in controlled smoke-test mode, not full
dataset generation mode.

% =========================================================
% 6. Notebook Pipeline
% =========================================================
\section{Notebook Pipeline}

The current notebook flow is:

\begin{enumerate}
    \item Install dependencies.
    \item Check CUDA runtime.
    \item Login to Hugging Face if needed.
    \item Resolve CityPersons dataset root.
    \item Define configuration and prompt templates.
    \item Scan CityPersons image files and labels.
    \item Build image records.
    \item Build augmentation jobs.
    \item Split jobs across available GPUs.
    \item Load Stable Diffusion 3.5 pipeline.
    \item Load YOLOv8m-seg model.
    \item For each image and variant, select a smart insertion region.
    \item Create a context crop around the target region.
    \item Generate a pedestrian using SD3.5.
    \item Run YOLOv8m-seg on the generated crop.
    \item Select the person mask closest to the target insertion box.
    \item If the person mask is invalid, retry generation.
    \item Color-match the generated person to the local scene.
    \item Paste only the person pixels onto the original image.
    \item Add contact shadow if enabled.
    \item Save augmented image, comparison pair, debug strip, and manifest entry.
\end{enumerate}

\begin{figure}[H]
\centering
\begin{tabular}{|p{0.9\linewidth}|}
\hline
\textbf{Current Notebook Pipeline} \\[4pt]
CityPersons Image \\
$\downarrow$ \\
Smart Placement \\
$\downarrow$ \\
Context Crop Construction \\
$\downarrow$ \\
SD3.5 Img2Img / Inpainting Generation \\
$\downarrow$ \\
YOLOv8m-seg Person Segmentation \\
$\downarrow$ \\
Person Mask Validation and Retry \\
$\downarrow$ \\
Color Matching \\
$\downarrow$ \\
Person-only Composite onto Original Image \\
$\downarrow$ \\
Save Image, Comparison Pair, Patch Debug, and Manifest \\
\hline
\end{tabular}
\caption{Current pipeline implemented in the notebook.}
\label{fig:current_pipeline}
\end{figure}

% =========================================================
% 7. Pipeline Evolution
% =========================================================
\section{Pipeline Evolution}

The current notebook is the result of multiple design iterations. Each version
was created to solve a specific failure mode.

\subsection{Full-Image Img2Img}

The earliest approach used full-image image-to-image generation. The model saw
the entire scene and was prompted to add a pedestrian.

\subsubsection{Problem}

At low strength, the model often did not add a visible person. At high strength,
the model added stronger pedestrian content but modified background regions.
Road texture, buildings, vehicles, and scene layout could change.

\subsubsection{Lesson}

Full-image diffusion is too unconstrained for background-preserving pedestrian
insertion. Prompting ``preserve background'' is not enough if the model is
allowed to denoise the entire image.

\subsection{Patch-Based Generation}

The next approach generated only a local patch. This preserved the background
outside the patch.

\subsubsection{Problem}

The background inside the patch could still be changed. If the patch was too
small, the generated person became weak or incomplete. If the patch was too
large, background corruption increased.

\subsubsection{Lesson}

Patch-based generation reduces the affected area but does not solve the
fundamental problem of generated background replacing original background.

\subsection{Human-Mask Inpainting}

A tighter human-shaped mask was tested to protect the background.

\subsubsection{Problem}

The mask was too restrictive. The model often generated ghost-like or
skeleton-like pedestrians. Some outputs showed only partial limbs or faint
silhouettes.

\subsubsection{Lesson}

Stable Diffusion needs spatial freedom to generate a complete body. A very tight
mask protects the background but harms person generation.

\subsection{Bounding-Box Inpainting}

The human-shaped mask was replaced by a rounded bounding-box mask.

\subsubsection{Improvement}

The pedestrian became clearer because the model had more room to generate a
complete person.

\subsubsection{Remaining Problem}

Background pixels inside the box could still be rewritten. Directly pasting the
generated box into the original image could create local artifacts.

\subsection{Flexible Scale Envelope}

The bounding box was then relaxed into a flexible scale envelope. Instead of
forcing the model to fill an exact person box, the notebook gives the model a
larger approximate region.

The current notebook uses:

\begin{lstlisting}[language=Python,caption={Current scale envelope values}]
SCALE_ENVELOPE_HEIGHT_MULT = 1.16
SCALE_ENVELOPE_WIDTH_MULT = 1.30
\end{lstlisting}

\subsubsection{Lesson}

The bbox should guide the model, not over-constrain it. Flexible envelopes allow
the model to generate more natural pose and body shape.

\subsection{Context Person Composite}

The final design is context person composite. In this design, SD3.5 is allowed
to generate inside a context crop, but the final image only uses the segmented
person pixels from the generated crop.

This separates:

\begin{itemize}
    \item person generation, handled by SD3.5;
    \item person extraction, handled by YOLOv8m-seg;
    \item background preservation, handled by original-image compositing.
\end{itemize}

% =========================================================
% 8. Smart Placement
% =========================================================
\section{Smart Placement}

Smart Placement selects a plausible insertion position before generation. This
is necessary because random insertion may place a pedestrian on a car, in the
sky, inside a building, or at an impossible scale.

\subsection{Perspective-Aware Scale}

The current notebook uses perspective priors:

\begin{lstlisting}[language=Python,caption={Current perspective scale values}]
PERSPECTIVE_SCALE_FAR = 0.46
PERSPECTIVE_SCALE_NEAR = 1.08
\end{lstlisting}

The scale is computed from the pedestrian ground point:

\begin{equation}
t =
\frac{y - y_{\text{far}}}
{y_{\text{near}} - y_{\text{far}}}
\end{equation}

\begin{equation}
s(y) =
s_{\text{far}} + t(s_{\text{near}} - s_{\text{far}})
\end{equation}

where $y$ is the vertical ground position of the inserted pedestrian.

\subsection{Variant-Specific Placement}

The notebook supports multiple variants:

\begin{itemize}
    \item \texttt{add\_single\_pedestrian};
    \item \texttt{add\_two\_pedestrians};
    \item \texttt{add\_small\_group};
    \item \texttt{add\_occluded\_pedestrian};
    \item \texttt{add\_distant\_pedestrian};
    \item \texttt{add\_near\_pedestrian}.
\end{itemize}

However, the current smoke test only uses:

\begin{itemize}
    \item \texttt{add\_single\_pedestrian};
    \item \texttt{add\_near\_pedestrian}.
\end{itemize}

This makes sense because near-person and single-person outputs are easier to
inspect during debugging.

% =========================================================
% 9. Prompt Engineering
% =========================================================
\section{Prompt Engineering}

Prompt design is important because SD3.5 has a token limit and can ignore weak
instructions. Earlier prompts were too long and repeated background preservation
phrases. The current pipeline reduces the need for long preservation prompts
because the final background is taken from the original image.

\subsection{Prompt Design Principle}

The prompt should focus on generating the new pedestrian:

\begin{itemize}
    \item realistic person;
    \item full body;
    \item grounded feet;
    \item correct scale;
    \item visible body parts;
    \item compatible with road or sidewalk context.
\end{itemize}

The pipeline, not the prompt alone, handles background preservation.

\subsection{Negative Prompt}

The negative prompt targets observed failure modes:

\begin{itemize}
    \item empty mask;
    \item unchanged image;
    \item invisible person;
    \item ghost person;
    \item cropped body;
    \item missing head;
    \item missing arms;
    \item missing legs;
    \item missing feet;
    \item floating person;
    \item wrong scale;
    \item bad anatomy;
    \item background change.
\end{itemize}

% =========================================================
% 10. YOLOv8m-seg Person Extraction
% =========================================================
\section{YOLOv8m-seg Person Extraction}

\subsection{Why YOLOv8m-seg Is Used}

The current notebook uses YOLOv8m-seg, not YOLOv8n-seg. The reason is that mask
quality is important in this pipeline. If the mask is poor, the pasted person
may have halos, missing parts, or background leakage.

YOLOv8m-seg provides stronger segmentation than the nano variant, while still
being practical in the notebook environment.

\subsection{Confidence and Mask Threshold}

The current thresholds are:

\begin{lstlisting}[language=Python,caption={Current YOLO segmentation thresholds}]
CONTEXT_PERSON_MIN_CONFIDENCE = 0.12
CONTEXT_PERSON_MASK_THRESHOLD = 0.32
\end{lstlisting}

The confidence threshold is relatively low because generated pedestrians may be
imperfect, especially during early attempts. A low threshold helps recover
usable person masks, while the retry and validation logic prevents weak masks
from being accepted silently.

\subsection{Person Selection}

YOLO may detect multiple people in the generated crop. The notebook selects the
person closest to the intended insertion region. This prevents existing people
or unrelated detections from being selected.

Conceptually:

\begin{lstlisting}[language=Python,caption={Person selection concept}]
def select_generated_person_mask(generated_crop, target_bbox):
    detections = yolo_segment(generated_crop)

    best_mask = None
    best_score = -1

    for detection in detections:
        if detection.class_name != "person":
            continue

        overlap_score = iou(detection.bbox, target_bbox)
        distance_score = center_similarity(detection.bbox, target_bbox)
        score = overlap_score + distance_score + detection.confidence

        if score > best_score:
            best_score = score
            best_mask = detection.mask

    return best_mask
\end{lstlisting}

% =========================================================
% 11. Person-Only Compositing
% =========================================================
\section{Person-Only Compositing}

\subsection{Core Idea}

The generated crop is not directly pasted into the original image. Instead, the
pipeline uses the YOLOv8m-seg mask to paste only the generated pedestrian pixels.

This design is the main reason the background can be preserved.

\begin{figure}[H]
\centering
\begin{tabular}{|p{0.9\linewidth}|}
\hline
\textbf{Person-only Composite Logic} \\[4pt]
Original image provides the final background. \\
Generated crop provides the candidate pedestrian. \\
YOLOv8m-seg extracts the person mask. \\
Only person pixels are pasted onto the original image. \\
Everything outside the person mask remains from the original image. \\
\hline
\end{tabular}
\caption{Person-only compositing strategy.}
\end{figure}

\subsection{Color Matching}

The notebook enables color matching:

\begin{lstlisting}[language=Python,caption={Current color matching flag}]
COLOR_MATCH_PERSON_TO_SCENE = True
\end{lstlisting}

This step adjusts the generated person's color statistics toward the local
source scene. It helps reduce the mismatch between generated clothing/person
pixels and the surrounding lighting.

\subsection{Contact and Grounding}

The prompt and placement logic encourage the pedestrian's feet to touch the
ground. If the generated person looks like a sticker, contact shadow can improve
integration. However, shadow should be light and local. Strong blending can
damage detector usefulness.

% =========================================================
% 12. Context Crop Design
% =========================================================
\section{Context Crop Design}

The current notebook uses a relatively large context crop:

\begin{lstlisting}[language=Python,caption={Current context crop settings}]
CONTEXT_CROP_EXPAND = 3.4
CONTEXT_INPAINT_MASK_PADDING = 46
\end{lstlisting}

\subsection{Why Context Crop Is Large}

A larger crop allows SD3.5 to see the surrounding road, sidewalk, vehicle scale,
and scene perspective. This is important because the model needs context to
generate a pedestrian that fits the scene.

\subsection{Trade-Off}

A larger crop increases context understanding but also gives the model more
background information to modify internally. However, because the final image
uses person-only compositing, internal crop background changes are less harmful.
The generated background is not directly trusted.

\subsection{Current Design Rationale}

The current configuration uses large context but controlled compositing:

\begin{itemize}
    \item large crop for better generation;
    \item YOLOv8m-seg for person extraction;
    \item original image for final background;
    \item retry if no valid person is detected.
\end{itemize}

% =========================================================
% 13. Retry Logic
% =========================================================
\section{Retry Logic}

The current notebook uses:

\begin{lstlisting}[language=Python,caption={Current retry configuration}]
CONTEXT_GENERATION_RETRIES = 4
\end{lstlisting}

This means the notebook can retry generation when no valid person mask is found.
Retry is necessary because diffusion generation is stochastic. Some attempts may
produce no person, a weak person, or a person that YOLO cannot segment.

\subsection{Retry Purpose}

Retry logic prevents the notebook from accepting bad samples too easily.
Instead of saving an unchanged or invalid output, the notebook attempts another
generation.

\subsection{Failure Cases That Trigger Retry}

Retry can be triggered by:

\begin{itemize}
    \item no YOLO person mask;
    \item person mask too small;
    \item person too faint;
    \item generated person far from target bbox;
    \item segmentation failure.
\end{itemize}

% =========================================================
% 14. Debug and Output
% =========================================================
\section{Debug and Output}

The current notebook enables patch debug:

\begin{lstlisting}[language=Python,caption={Current debug settings}]
SAVE_PATCH_DEBUG = True
PATCH_DEBUG_MAX_ITEMS = 24
\end{lstlisting}

\subsection{Why Debug Output Is Important}

Debug output is essential because the final augmented image alone does not show
where the failure occurred. A failure may come from:

\begin{itemize}
    \item Smart Placement;
    \item SD3.5 generation;
    \item YOLO segmentation;
    \item mask selection;
    \item color matching;
    \item compositing.
\end{itemize}

Patch debug strips make it possible to inspect intermediate steps such as source
crop, generated crop, person mask, and final composite.

\subsection{Manifest Logging}

The notebook writes manifest information so that each generated sample can be
traced back to its source image, variant, seed, generation parameters, and output
path.

Recommended manifest fields include:

\begin{lstlisting}[language={},caption={Recommended manifest columns}]
source_path
output_path
comparison_path
variant
strength
guidance_scale
num_inference_steps
seed
target_bbox
selected_person_bbox
detected_person_confidence
person_mask_area_ratio
context_retry_count
segmentation_success
fallback_used
accepted
reject_reason
\end{lstlisting}

% =========================================================
% 15. Failure Analysis
% =========================================================
\section{Failure Analysis}

\subsection{Background Corruption}

\textbf{Symptom:} Road, buildings, vehicles, or sidewalk textures are modified.

\textbf{Cause:} Full-image or patch-level generation allows diffusion to rewrite
background pixels.

\textbf{Current Fix:} Use \texttt{context\_person\_composite}. The generated
crop is used only to extract the person; the original image remains the
background.

\subsection{Ghost Pedestrian}

\textbf{Symptom:} The generated pedestrian appears as a faint silhouette.

\textbf{Cause:} Generation signal is too weak, mask is too tight, or prompt is
not directive enough.

\textbf{Current Fix:} Use context crop, flexible scale envelope, stronger
person-focused prompt, and retry logic.

\subsection{Partial Body}

\textbf{Symptom:} Missing head, missing legs, missing arms, or cropped body.

\textbf{Cause:} Inpaint region too small, generation region too constrained, or
segmentation mask removes body parts.

\textbf{Current Fix:} Use context crop expansion, mask padding, and full-body
prompt. The current mask padding is \texttt{46}, giving the model more space.

\subsection{Wrong Scale}

\textbf{Symptom:} Pedestrian is too small, too large, or inconsistent with
perspective.

\textbf{Cause:} Placement and scale estimation do not match scene geometry.

\textbf{Current Fix:} Use perspective scale priors with
\texttt{PERSPECTIVE\_SCALE\_FAR=0.46} and
\texttt{PERSPECTIVE\_SCALE\_NEAR=1.08}, plus flexible scale envelope.

\subsection{Segmentation Failure}

\textbf{Symptom:} YOLO does not detect the generated pedestrian.

\textbf{Cause:} The person is too faint, too small, incomplete, or confidence is
below threshold.

\textbf{Current Fix:} Use YOLOv8m-seg, lower confidence threshold
\texttt{0.12}, mask threshold \texttt{0.32}, and retry up to four times.

\subsection{Prompt Token Limit}

\textbf{Symptom:} Token sequence length warning.

\textbf{Cause:} Prompt too long.

\textbf{Fix:} Keep prompts short and focus on the object. Let the pipeline
handle background preservation.

\subsection{Large Output Folder}

\textbf{Symptom:} Output zip becomes too large.

\textbf{Cause:} Zipping all of \texttt{/kaggle/working} instead of only the
augmentation output directory.

\textbf{Fix:} Zip only the configured \texttt{OUTPUT\_DIR}.

% =========================================================
% 16. Evaluation Strategy
% =========================================================
\section{Evaluation Strategy}

\subsection{Why Evaluation Is Needed}

Generated images should not be accepted only by visual inspection. For Camera AI
augmentation, the output must be useful for detection and tracking. A visually
interesting image may still be useless if the detector cannot detect the
inserted person or if the background changes too much.

\subsection{Proposed Metrics}

\begin{table}[H]
\centering
\caption{Evaluation metrics for the current pipeline.}
\label{tab:evaluation_metrics}
\begin{tabular}{lll}
\toprule
\textbf{Metric} & \textbf{Tool} & \textbf{Purpose} \\
\midrule
Person detection rate & YOLO detector & Check if person is visible \\
Person segmentation confidence & YOLOv8m-seg & Check mask reliability \\
Person mask area ratio & YOLOv8m-seg & Reject ghost/tiny persons \\
Background SSIM & Image similarity & Check background preservation \\
Background LPIPS & Perceptual metric & Check perceptual drift \\
Retry count & Notebook log & Check generation stability \\
Fallback rate & Manifest & Check failure frequency \\
Variant correctness & YOLO + manual review & Check if variant matches prompt \\
Downstream mAP & Detector retraining & Final usefulness \\
\bottomrule
\end{tabular}
\end{table}

\subsection{Smoke-Test Evaluation}

The current notebook is configured for a small smoke test:

\begin{itemize}
    \item \texttt{total\_images=10};
    \item variants: \texttt{add\_single\_pedestrian},
          \texttt{add\_near\_pedestrian};
    \item split: \texttt{train}.
\end{itemize}

The purpose of this test is not to produce the final dataset. It is to verify:

\begin{itemize}
    \item whether SD3.5 generates a visible person;
    \item whether YOLOv8m-seg can extract the person;
    \item whether person-only composite preserves background;
    \item whether scale and ground contact look reasonable;
    \item whether retry rate is acceptable.
\end{itemize}

% =========================================================
% 17. Discussion
% =========================================================
\section{Discussion}

\subsection{Why the Current Pipeline Is More Correct Than Full Img2Img}

Full-image img2img gives Stable Diffusion control over the entire image. This
makes it hard to preserve scene geometry. The current pipeline gives SD3.5
control only inside a context crop, then extracts the person and restores the
background from the original image.

This is more appropriate for dataset augmentation because the original scene
geometry is still reliable.

\subsection{Why YOLOv8m-seg Matters}

The current notebook depends on segmentation quality. If the segmentation mask
is poor, the final composite will be poor. This explains why the notebook uses
\texttt{yolov8m-seg.pt}. The medium segmentation model is a better choice than
the nano model when mask quality is important.

\subsection{Why Strength Is Lower in the Current Notebook}

Earlier versions used very high strength to force the model to generate a
person. The current notebook uses lower global strength:

\begin{lstlisting}[language=Python]
AUGMENTATION_STRENGTH = 0.72
GUIDANCE_SCALE = 7.2
\end{lstlisting}

The reason is that the current pipeline uses larger context, better segmentation,
and retry logic. It does not need to over-force generation as aggressively. The
variant-specific strengths still increase for harder cases such as near
pedestrians or groups.

\subsection{Why LoRA Is Still Optional}

LoRA may become useful if the current pipeline fails repeatedly. However, many
previous problems were not caused by lack of model knowledge. They were caused
by editing-control issues:

\begin{itemize}
    \item background rewrite;
    \item bad placement;
    \item wrong scale;
    \item weak segmentation;
    \item poor compositing;
    \item prompt length.
\end{itemize}

These problems are better solved by pipeline design than by immediate
fine-tuning.

% =========================================================
% 18. Current State
% =========================================================
\section{Current State}

The current notebook state is:

\begin{itemize}
    \item Main mode: \texttt{context\_person\_composite}.
    \item Main generation model: Stable Diffusion 3.5 Medium.
    \item Text encoder setting: \texttt{USE\_T5=False}.
    \item Resolution: \texttt{448}.
    \item Segmentation model: \texttt{yolov8m-seg.pt}.
    \item Context crop expansion: \texttt{3.4}.
    \item Mask padding: \texttt{46}.
    \item Retry count: \texttt{4}.
    \item Color matching: enabled.
    \item Patch debug: enabled.
    \item Smoke test: 10 images.
    \item Current variants: \texttt{add\_single\_pedestrian} and
          \texttt{add\_near\_pedestrian}.
\end{itemize}

This means the notebook is currently in a controlled evaluation stage. It is not
yet running full-scale augmentation.

% =========================================================
% 19. Future Work
% =========================================================
\section{Future Work}

\subsection{Run Larger Evaluation}

The next step should be progressive scaling:

\begin{enumerate}
    \item Run 10 images for smoke test.
    \item Review patch debug and comparison pairs.
    \item Run 24 or 40 images with more variants.
    \item Compute retry rate and segmentation success rate.
    \item Run 200--500 images only after quality is stable.
\end{enumerate}

\subsection{Improve Manifest Metrics}

The manifest should include more diagnostic fields:

\begin{itemize}
    \item retry count;
    \item segmentation success;
    \item selected person confidence;
    \item person mask area ratio;
    \item target bbox;
    \item selected person bbox;
    \item reject reason;
    \item fallback flag;
    \item background similarity metrics.
\end{itemize}

\subsection{Downstream Detector Training}

After enough accepted samples are generated, the next experiment should compare:

\begin{itemize}
    \item detector trained without synthetic augmentation;
    \item detector trained with accepted synthetic images;
    \item validation performance on CityPersons;
    \item possible tracking evaluation on MOT17.
\end{itemize}

\subsection{LoRA Decision}

LoRA should only be considered if the current no-LoRA baseline fails
systematically. Possible triggers include:

\begin{itemize}
    \item low person generation success rate;
    \item high retry rate;
    \item repeated ghost pedestrians;
    \item repeated missing body parts;
    \item poor realism even when placement and compositing are correct.
\end{itemize}

% =========================================================
% 20. Conclusion
% =========================================================
\section{Conclusion}

This project develops a synthetic pedestrian augmentation pipeline for
CityPersons using Stable Diffusion 3.5 Medium. The current notebook implements a
context-aware generation and YOLOv8m-seg based person-only compositing pipeline,
rather than a LoRA-first training pipeline.

The main design is to let Stable Diffusion generate a pedestrian inside a large
context crop while using YOLOv8m-seg to extract the generated person. The final
image preserves the original background by pasting only person pixels back onto
the source frame. This addresses the main weakness of full-image img2img:
background corruption.

The current configuration uses \texttt{yolov8m-seg.pt}, \texttt{RESOLUTION=448},
\texttt{AUGMENTATION\_STRENGTH=0.72}, \texttt{GUIDANCE\_SCALE=7.2},
\texttt{CONTEXT\_CROP\_EXPAND=3.4}, \texttt{CONTEXT\_INPAINT\_MASK\_PADDING=46},
and \texttt{CONTEXT\_GENERATION\_RETRIES=4}. The notebook is currently configured
to run a 10-image smoke test on the training split with
\texttt{add\_single\_pedestrian} and \texttt{add\_near\_pedestrian} variants.

The most important conclusion is that background preservation is an architectural
problem, not only a prompt-engineering problem. The second conclusion is that
LoRA fine-tuning is optional and should only be introduced after the current
baseline is measured and shown to fail systematically.

% =========================================================
% References
% =========================================================
\begin{thebibliography}{9}

\bibitem{citypersons}
S. Zhang, R. Benenson, and B. Schiele,
``CityPersons: A Diverse Dataset for Pedestrian Detection,''
in \textit{Proceedings of the IEEE Conference on Computer Vision and Pattern Recognition},
2017.

\bibitem{stablediffusion35}
Stability AI,
``Stable Diffusion 3.5 Medium,''
2024.

\bibitem{yolov8}
G. Jocher et al.,
``Ultralytics YOLOv8,''
2023.

\bibitem{lpips}
R. Zhang, P. Isola, A. A. Efros, E. Shechtman, and O. Wang,
``The Unreasonable Effectiveness of Deep Features as a Perceptual Metric,''
in \textit{Proceedings of the IEEE Conference on Computer Vision and Pattern Recognition},
2018.

\bibitem{mot17}
A. Milan, L. Leal-Taixe, I. Reid, S. Roth, and K. Schindler,
``MOT16: A Benchmark for Multi-Object Tracking,''
arXiv preprint arXiv:1603.00831, 2016.

\end{thebibliography}

\end{document}